\section{Frey representations in solvable quotients}
\subsection{A prosolvable Frey-three criterion}
\label{subsec:frey-three}
Let $n\ge3$ be odd. Identify $X_n$ with
$F_n:X^n+Y^n=Z^n$ and put $U=F_n-\{XYZ=0\}$.
For a point-or-cusp Selmer section $s$ of the geometrically
prosolvable quotient, use the Legendre family
$E_\lambda:y^2=x(x-1)(x-\lambda)$, where $\lambda=(X/Z)^n$.
All local statements below concern representatives of this section;
no lift to the full arithmetic fundamental group is assumed.
Write $\Pi_U^{\mathrm{sol}}$ for its arithmetic quotient.
\par\medskip\noindent\textit{The representation on the stated quotient.}\quad 

The geometric determinant of the Tate-module representation is trivial.
Every finite geometric image on \(E[3^m]\) is therefore a subgroup of
\(\mathrm{SL}_2(\mathbb Z/3^m)\). This group is solvable: its congruence
kernel over \(\mathbb F_3\) is a finite 3-group, while
\(\mathrm{SL}_2(\mathbb F_3)\) has a normal quaternion subgroup with
cyclic quotient and is solvable. Thus the representation factors through
\(\Pi_U^{\mathrm{sol}}\) at every level, and hence on \(T_3E\).

A quotient section \(s\) consequently supplies

\[
\rho_s:G_{\mathbb Q}\longrightarrow
\mathrm{GL}_2(\mathbb Z_3),\qquad \det\rho_s=\chi_3.
\]

No section of the full arithmetic fundamental group is assumed or
constructed. The action on a finite cyclic-subgroup scheme below is
likewise already defined on the prosolvable quotient.

\par\medskip\noindent\textit{The topology of solvable covers.}\quad 

\begin{lemma}[The full prosolvable topology on a solvable cover]\label{lem:solvable-open-topology} Let \(\Delta\) be a profinite group and
\(H\subseteq\Delta\) an open subgroup such that
\(\Delta/\operatorname{core}_\Delta(H)\) is solvable. If \(R_\Delta\)
and \(R_H\) are the respective kernels of the maximal prosolvable
quotients, then

\[
R_H=R_\Delta\subseteq H.
\]

In particular \(H^{\mathrm{sol}}\) is the open subgroup \(H/R_\Delta\)
of \(\Delta^{\mathrm{sol}}\).

\end{lemma}\begin{proof} The solvable permutation quotient first gives
\(R_\Delta\subseteq H\). Let \(K\triangleleft H\) be open with \(H/K\)
solvable. The core \(C=\bigcap_{\delta\in\Delta}\delta K\delta^{-1}\) is
a finite intersection because \(K\) is open and normal in \(H\). Set
\(L=\operatorname{core}_\Delta(H)\). Then

\[
L/C\hookrightarrow
\prod_{\delta\in\Delta/H}
 \delta H\delta^{-1}/\delta K\delta^{-1}.
\]

The right side is solvable, as is \(\Delta/L\). Hence \(\Delta/C\) is
solvable and \(R_\Delta\subseteq C\subseteq K\). Intersecting over \(K\)
proves \(R_\Delta\subseteq R_H\).

Conversely, for every open normal \(M\triangleleft\Delta\) with
\(\Delta/M\) solvable, the quotient \(H/(H\cap M)\) is solvable.
Therefore \(R_H\subseteq\bigcap_M(H\cap M)=R_\Delta\). This proves the
lemma.

\end{proof}

For a finite cover with solvable geometric Galois closure, the lemma
applies to geometric fundamental groups. Quotienting its full arithmetic
fundamental group by this common kernel gives exactly the corresponding
open subgroup of the base's geometrically prosolvable arithmetic group.
This is an identification of entire prosolvable quotients, not just
their particular finite quotients.

\par\medskip\noindent\textit{A finite sufficient condition.}\quad 

\begin{proposition}[A stable cyclic subgroup modulo nine]\label{prop:frey-three-mod-nine} If the representation \(\rho_s\bmod9\) has a
\(G_{\mathbb Q}\)-stable cyclic subgroup of order 9, then \(s\) is
globally cuspidal.
\end{proposition}

The hypothesis is equivalent to a stable rank-one direct summand of
\((\mathbb Z/9)^2\). It is stronger than a stable line modulo 3. For
example, the inverse image of a Borel modulo 3 in
\(\mathrm{GL}_2(\mathbb Z/9)\) preserves the residual line, but the
congruence element \(1+3E_{21}\) moves each of its lifts. This is a
finite arithmetic sufficient condition, not a reduction of the
difficulty of proving mod-3 reducibility.

\begin{proof}\par\medskip\noindent\textit{The correct neighborhood.}\quad 

Let \(\mathcal S\) be the finite étale scheme over \(U\) parametrizing
cyclic order-9 subgroups of \(E[9]\). Choose the geometric point \(C\)
fixed by the pulled-back action of \(s(G_{\mathbb Q})\). Its stabilizer
is an open subgroup of \(\Pi_U^{\mathrm{sol}}\) containing
\(s(G_{\mathbb Q})\).

Take the corresponding component \(V\to U\). It is defined over
\(\mathbb Q\) and geometrically connected: the arithmetic orbit of \(C\)
is its geometric orbit, since the arithmetic quotient is generated by
its geometric subgroup together with the section and the latter fixes
\(C\). Equivalently, the inverse-image stabilizer in \(\pi_1(U)\)
surjects onto \(G_{\mathbb Q}\).

The geometric permutation group of \(\mathcal S\) is a quotient of a
subgroup of \(\mathrm{SL}_2(\mathbb Z/9)\), so its Galois closure is
solvable. Lemma~\ref{lem:solvable-open-topology} identifies the open arithmetic stabilizer with
\(\Pi_V^{\mathrm{sol}}\). The given section therefore lifts to a section
of that precise group. A stable subgroup in the pulled-back
representation has thus been converted into a compatible lift to the
level curve; this conversion is justified by a stabilizer argument, not
by assuming a full-section lift.

The lifted section is locally point-or-cusp geometric. For an interior
local point, fixedness of the corresponding finite fibre gives a
rational point of the cover. At a boundary point, the same statement
uses the finite set of prolongations and their stabilizers. One can also
lift the local cusp section before passing to the finite cover: cusp
inertia injects into the prosolvable quotient here, since each puncture
has a nonzero primitive class in the geometric abelianization and the
curve has more than one puncture. Its intersection with the stabilizer
projects onto the local Galois group and therefore gives a rational
boundary lift. This argument does not require that the entire cusp
decomposition group be contained in the neighborhood.

\par\medskip\noindent\textit{The map to the modular curve.}\quad 

Over \(V\), the elliptic curve has its original full rational 2-torsion
and the tautological cyclic subgroup \(C_9\). Choose independent nonzero
2-torsion sections \(P,Q\). Let

\[
\phi:E\longrightarrow E'=E/\langle P\rangle
\]

be the 2-isogeny, with dual \(\widehat\phi\). The subgroup

\[
C_4=\widehat\phi^{-1}(\langle Q\rangle)
\]

is cyclic of order 4. Indeed, for \(\widehat\phi(R)=Q\) one has
\(2R=\phi(Q)\ne0\), so \(R\) has order 4. The image \(C'_9=\phi(C_9)\)
is cyclic of order 9. Their sum is a cyclic subgroup \(C_{36}\) of order
36 because their orders are coprime.

This gives an algebraic map \(V\to Y_0(36)\), and thus a map from the
smooth proper compactification \(\overline V\) to \(X_0(36)\). It is
nonconstant. If \(j(E')\) were constant, the degree-2 modular relation
would force \(j(E)\) algebraic over the constant field, contrary to the
nonconstant Legendre parameter.

The modular curve has genus one: the standard index is 72, the number of
cusps is 12, and both elliptic-point contributions vanish, giving
\(1+72/12-12/2=1\). \citet[Corollary~8.8, p.~385]{Stoll2007} proves
that every positive-genus \(X_0(N)\) over \(\mathbb Q\) is excellent for
finite abelian descent. Proposition~8.5, p.~384, of the same source consequently
applies to every proper curve mapping nonconstantly to \(X_0(36)\). Its
real-place proof correction is in \citet[\S3]{StollErrata2017} and does not
alter the assertion.

\par\medskip\noindent\textit{Cusp containment.}\quad 

Let \(W\to V\) run through finite neighborhoods of the lifted
prosolvable section. One can use a cofinal family of open characteristic
subgroups of the finitely generated geometric prosolvable group and
adjoin the section image. Their geometric Galois closures are solvable.
Lemma~\ref{lem:solvable-open-topology} applies again, so each carries its induced entire prosolvable
section. The compactification \(\overline W\) maps nonconstantly to
\(X_0(36)\); hence

\[
\overline W(\mathbb A_{\mathbb Q})_\bullet^{f\text{-ab}}
=\overline W(\mathbb Q).
\]

The adelic point attached to the lifted local sections survives finite
abelian descent. This assertion only uses the given prosolvable section:
any finite abelian torsor on \(\overline W\), pulled back to \(W\),
factors through that quotient, and the section supplies its common
global twist with local lifts. No non-solvable torsor or full
fundamental-group section is required.

Consequently every \(\overline W\) has a rational point compatible with
the chosen local section. To make compatibility explicit, fix one finite
place and a point or cusp representative of the local section of \(V\).
Its lifts through all section neighborhoods give compatible local
points. Excellence gives rational points with those localizations;
injectivity of rational-point localization then makes the rational
points compatible under every transition map.

Fix their image \(z\in\overline V(\mathbb Q)\). Above \(z\), the
compactified finite geometric fibres form a projective system of finite
sets. A rational lift is a section-fixed orbit at the corresponding
level. The compatible rational points thus determine a fixed
prolongation in the inverse-limit fibre. Its stabilizer is a
decomposition group. Hence the lifted section image is contained in that
decomposition group.

The direction is

\[
s(G_{\mathbb Q})\subseteq D_{\widetilde z},
\]

not the reverse; cusp inertia need not lie in a section image. This
finite-orbit argument remains valid at ramified boundary fibres.

Finally the image of \(z\) on \(F_n\) is a rational coordinate cusp,
by the classical Fermat theorem \citep[Theorem~0.5, p.~448]{Wiles1995}. Thus \(z\) belongs to the boundary of \(V\), and
its image section on \(U\) is cuspidal. This proves the proposition.\end{proof}

\begin{corollary}\label{cor:frey-three-irreducible} If \(\rho_s\) is reducible over \(\mathbb Q_3\),
then \(s\) is cuspidal.

\end{corollary}\begin{proof}
Indeed, an invariant \(\mathbb Q_3\)-line intersects the stable Tate
lattice in a saturated rank-one \(\mathbb Z_3\)-submodule. Its reduction
modulo 9 is a stable direct summand, so Proposition~\ref{prop:frey-three-mod-nine} applies. Therefore
every hypothetical noncuspidal section has \(\mathbb Q_3\)-irreducible
\(\rho_s\).
\end{proof}

\par\medskip\noindent\textit{Density-one integrality.}\quad 

The next argument explicitly handles disconnected algebraic monodromy
and representations that are irreducible over \(\mathbb Q_3\) but
reducible over its algebraic closure.

\begin{lemma}[The squared-trace zero set]\label{lem:frey-three-trace-density} Let \(H\) be a compact subgroup of
\(\mathrm{GL}_2(\mathbb Q_3)\), acting irreducibly over \(\mathbb Q_3\),
and suppose \(\det H\) is infinite. Then

\[
T=\{h\in H:\operatorname{tr}(h)^2=(1+\det h)^2\}
\]

has Haar measure zero.

\end{lemma}\begin{proof} Let \(G\) be the algebraic Zariski closure of \(H\), and
\(G^\circ\) its connected identity component. Its determinant is
\(\mathbb G_m\): the determinant of \(G\) is infinite, and a finite
component group cannot carry that dimension.

The unipotent radical of \(G^\circ\) is trivial. Otherwise a nontrivial
unipotent subgroup in dimension two has a unique fixed line, which is
preserved by all of \(G\), since the radical is characteristic in
\(G^\circ\). That line descends to \(\mathbb Q_3\) and contradicts
irreducibility.

If the connected reductive group \(G^\circ\) is not a torus, its derived
group contains \(\mathrm{SL}_2\). On every coset \(gG^\circ\), the trace
varies along \(g\mathrm{SL}_2\) while the determinant is constant.
Therefore the polynomial

\[
F(g)=\operatorname{tr}(g)^2-(1+\det g)^2
\]

cannot vanish identically on that component.

It remains to treat a torus, after extension to
\(\overline{\mathbb Q}_3\). If it acts by scalars, the scalar parameter
\(t\) gives

\[
F(tg)=t^2\operatorname{tr}(g)^2-(1+t^2\det g)^2,
\]

a nonzero polynomial. If it has two distinct weights, each \(g\in G\)
either preserves or exchanges their weight lines. In an exchanging
component the trace is zero, whereas the determinant varies, so \(F\)
again is not identically zero. In a preserving component let the
eigenvalues be \(a\alpha(t)\) and \(b\beta(t)\). Then

\[
F(gt)=-
\bigl(a^2\alpha(t)^2-1\bigr)
\bigl(b^2\beta(t)^2-1\bigr).
\]

Identical vanishing on the connected torus would force one weight to be
trivial. Since the determinant varies, precisely one weight is trivial.
Its weight space is the intrinsic invariant subspace of \(G^\circ\); it
is preserved by all of \(G\) and defined over \(\mathbb Q_3\). This
contradicts irreducibility. This last observation also covers a nonsplit
torus and the absolutely reducible case; no absolute-irreducibility
assumption has been made.

Thus \(F\) is nonzero on every algebraic component. Every open subgroup
of \(H\) has Zariski closure containing \(G^\circ\), since it has finite
index. Hence \(F\) cannot vanish identically on an analytic open subset
of any component of \(H\). A compact linear 3-adic group is a 3-adic
analytic group, and a nonzero analytic function on a chart has a zero
locus of Haar measure zero. The latter fact follows by induction on the
chart dimension, applying the one-variable zero theorem to fibres and
then Fubini. A finite chart cover proves the result. \citet[\S2.4, author pp.~4--5]{KhareRajan2001} use the same squared-trace hypersurface and
finite-level density mechanism in the open-image case; the component
analysis above supplies the extra cases needed here.

\end{proof}

\begin{proposition}[Density-one integrality]\label{prop:frey-three-density} If \(\rho_s\) is irreducible over
\(\mathbb Q_3\), the set of primes at which its associated proper local
point is cuspidal or reduces to the boundary of \(U\) has upper
Dirichlet density zero.

\end{proposition}\begin{proof} Discard the finite set of primes where the chosen
curve/family model is bad, together with 3. At a boundary or degenerate
place \(q\), the local Tate representation, possibly with a quadratic
twist, is upper triangular with diagonal characters
\(\delta_q\chi_3,\delta_q\). Thus every arithmetic Frobenius lift has
squared trace \((1+q)^2\). The cusp calculation has the same diagonal
form. These are the local calculations in \citet[\S4.3, author pp.~7--8]{Stix2015};
their validity is local and does not require a full global section. At the Legendre cusp at infinity, the quadratic twist may be ramified. No unramifiedness of that character is required: its square is one, so the squared-trace identity is unchanged, and the finite-level argument below allows its ramification.

For each \(m\), the representation modulo \(3^m\) has finite image and
is ramified at only finitely many primes. At a degenerate place outside
this finite exceptional set its Frobenius belongs to

\[
T_m=\{h\in \rho_s(G_{\mathbb Q}) :
 F(h)\equiv0\pmod {3^m}\}.
\]

Finite-level Chebotarev bounds the upper density of all degenerate
places by the Haar measure of \(T_m\). The closed sets \(T_m\) decrease
to the set \(T\) of Lemma~\ref{lem:frey-three-trace-density}, whose Haar measure is zero. Letting \(m\)
increase proves the proposition.
\end{proof}

The finite exceptional set is allowed to depend on \(m\). No finite
exceptional set for the entire 3-adic representation is used. Combining
Propositions~\ref{prop:frey-three-mod-nine} and~\ref{prop:frey-three-density} gives the proved reduction:

A noncuspidal prosolvable Selmer section has irreducible $\rho_s$
and is integral at a density-one set of primes. These are the modular
and density assertions used in Theorem~\ref{thm:prosolvable-frey-open-intro}.

Density-one integrality is not membership in
\(U(\mathbb A_{\mathbb Q})\), which would require integrality outside a
finite set.

\begin{remark}[The local parameter at seventeen]\label{rem:frey-three-seventeen-parameter}

Near the coordinate cusp \(\lambda=0\), put \(t=X/Z\). The Legendre
\(j\)-invariant and its Tate parameter give

\[
j=256\frac{(1-\lambda+\lambda^2)^3}
 {\lambda^2(1-\lambda)^2},
\qquad
q_E=2^{-8}t^{34}(1+O(t^{17})).
\]

Consequently, up to the fixed quadratic twist, the 3-adic extension
parameter for a tangential section is

\[
\kappa_3(2^{-8})+34\alpha_3.
\]

The coefficient 34 is a unit in \(\mathbb Z_3\). Thus the Fermat
exponent does not annihilate residual 3-adic tame monodromy. In the
usual exponent-prime Frey argument the corresponding divisibility kills
the residual inertia; that step has no analogue here.

At \(\mathbb Q_{17}\), the pro-3 completion of the unit group is zero,
since its finite part has order 16 and its principal-unit group is
pro-17. Thus \(H^1(\mathbb Q_{17},\mathbb Z_3(1))\cong\mathbb Z_3\) by
the valuation, and the constant term vanishes. Multiplication by 34 is
an automorphism. Local cuspidality at 17 therefore permits every 3-adic
Tate extension parameter; it gives a local invariant line, with no
global invariant line and no vanishing of local monodromy forced by this
calculation.

There is a concrete illustration of this distinction at the
representation level. The rational Legendre curve with \(\lambda=17\)
has full rational 2-torsion and split multiplicative reduction at 17.
Its 3-adic local representation there occurs in a cusp packet by the
preceding surjectivity. Yet its mod-3 representation has no global
stable line: its 3-division polynomial is

\[
3x^4-72x^3+102x^2-289.
\]

Modulo 5 its values at \(0,1,2,3,4\) are \(1,4,1,3,3\), so it has no
rational root, hence no rational cyclic 3-isogeny. This is not claimed
to arise from a Selmer section on \(U_{17}\); it only refutes the
implication from the local Tate shape at 17, determinant, and full
2-torsion alone to global residual reducibility.

Continuity gives finite ramification separately for each \(3^m\)-torsion
representation. It does not give a common finite ramification set for
the inverse limit. \citet[Theorem~1]{KhareRajan2001} prove density zero, rather than
finiteness, for ramification of semisimple \(p\)-adic representations.
\citet[Main Theorem and Remark~1, author p.~3]{KhareLarsenRamakrishna2005} construct, for coefficient primes at least 5,
infinitely ramified representations that are potentially semistable at
the coefficient prime and pure and rational at almost all unramified
primes. Their result is a counterexample to that general shortcut; its
stated prime restriction means it is not itself a 3-adic or
section-theoretic counterexample.

Nor does the finite-level modular curve associated only with a stable
line modulo 3 finish the argument. With full 2-torsion it gives level
12, whose modular curve has genus zero. The order-9 subgroup used above
gives level 36 and is a strictly stronger condition.

\end{remark}

\par\medskip\noindent\textit{The remaining sufficient condition.}\quad 

The sufficient conditions are global mod-9 Borel reduction, or
the stronger assumption of reducibility of the entire 3-adic
representation. Neither has been proved for an arbitrary prosolvable
Selmer section. The local theorem at seventeen by itself gives only a local line at
17. The next subsection proves that the remaining case has jointly open
two-adic and three-adic images. It has density-one integrality, but no
finite ramification set and no modularity theorem applicable solely from
the verified hypotheses.

This Frey-curve argument does not establish either sufficient condition
for an arbitrary prosolvable Selmer section. Exponent seventeen is
classified separately by the regular-prime theorem,
Theorem~\ref{thm:regular-solvable-intro}. Outside that theorem's range,
the simultaneous open-image case remains unresolved; no all-prime
finite-descent equality or global lifting theorem follows from the
Frey-curve argument alone.
